\chapter{Introduction}\label{ch:introduction}

\section{Superconducting quantum impurities}

A quantum dot, a molecule or a magnetic adatom coupled to superconductors is the
simplest system in which electron--electron interaction, superconducting pairing and
quantum coherence compete on an equal footing. The impurity level is described by the
Anderson model~\cite{anderson1961localized}: a single orbital that can hold zero, one or
two electrons, with charging energy $U$ for double occupancy. Coupled to a BCS
superconductor, the impurity binds Andreev or Yu--Shiba--Rusinov (YSR) states inside the
gap~\cite{yu1965bound,shiba1968classical,rusinov1969superconductivity,balatsky2006impurity}.
In the simplest spin-symmetric, single-orbital junction, the low-energy competition
is between a spin \emph{singlet} and a spin \emph{doublet}. A singlet can arise from
proximity pairing or from screening of a local moment; a doublet retains total
spin $1/2$ even when its dot-local polarization is reduced. Their crossing changes
the ground-state fermion parity. In a Josephson junction it is closely associated
with the change between 0- and $\pi$-junction behaviour, defined by the phase dependence
of the ground-state energy and current: the 0--$\pi$
transition~\cite{glazman1989resonant,spivak1991negative,rozhkov1999josephson,vecino2003josephson,choi2004kondo,siano2004josephson,bauer2007spectral,karrasch2008josephson,meng2009selfconsistent}.
Reviews of the theory and of the experiments are given in
Refs.~\cite{defranceschi2010hybrid,martinrodero2011josephson,meden2019andersonjosephson,prada2020andreev}.

Quantum-dot Josephson junctions have been realized in carbon nanotubes, semiconductor
nanowires and two-dimensional heterostructures, and the 0--$\pi$ transition, the
singlet--doublet level crossing and the subgap spectrum have been measured by transport,
tunnelling spectroscopy and circuit quantum electrodynamics
techniques~\cite{vandam2006supercurrent,cleuziou2006carbon,pillet2010andreev,deacon2010tunneling,lee2014spinresolved,delagrange2015manipulating,bargerbos2022singletdoublet}.
The same physics governs magnetic adatoms on superconducting surfaces probed by scanning
tunnelling microscopy~\cite{yazdani1997probing,franke2011competition,ruby2015tunneling,heinrich2018single}.

\section{Why drive the junction?}

The subgap states of a junction are the levels of a solid-state artificial atom, and
like the levels of an atom they are probed and controlled by oscillating fields.
Microwave radiation applied through the phase difference excites pairs of
quasiparticles into the Andreev levels of atomic contacts and
nanowires~\cite{bretheau2013exciting,bretheau2013supercurrent,janvier2015coherent,vanwoerkom2017microwave,hays2018direct}.
The even-parity pair transition is used in Andreev-level qubits; Andreev spin
qubits instead use two states within an odd-parity manifold, with spin-dependent
control matrix elements~\cite{desposito2001controlled,zazunov2003andreev,chtchelkatchev2003andreev,padurariu2010theoretical,hays2021coherent,pitavidal2023direct}.
Microwave probes also resolve the excitation spectra of interacting quantum-dot
junctions~\cite{tosi2019spinorbit,matutecanadas2022signatures,fatemi2022microwave,bargerbos2023spectroscopy,wesdorp2023dynamical}.
Strong periodic driving creates new states, the Floquet--Andreev states
observed in graphene junctions~\cite{park2022steady}, and microwave-assisted tunnelling
through YSR states has been observed with a scanning tunnelling
microscope~\cite{kot2020microwaveassisted,peters2020resonant,acerogonzalez2020photonassisted}.
Periodic driving is, more generally, a tool to engineer effective Hamiltonians, a
programme known as Floquet engineering~\cite{goldman2014periodically,bukov2015universal,eckardt2017colloquium,oka2019floquet,rudner2020band}.

These experiments raise several common questions. How are the many-body subgap states
dressed by the drive? Where do resonances occur and how large are the resulting
splittings? Which transitions are allowed by symmetry? How does the interaction modify
the answers obtained for non-interacting Andreev levels? Answering them requires the
spectrum of the \emph{periodically driven} many-body Hamiltonian, calculated by the
Floquet methods in \pkg{}.

\section{What the package computes}

\pkg{} treats a finite version of the model: each superconducting reservoir is
represented by a finite set of levels chosen to approximate its hybridization function. The
static solvers find low-lying many-body eigenstates in each symmetry sector, using either a
truncated quasiparticle (QP) expansion or the density-matrix renormalization group
(DMRG). Both solvers work with the same Hamiltonian and the same physical conventions.

When parameters are driven periodically, the code computes the \emph{quasienergies}
$\qe_\alpha$, the time-periodic \emph{Floquet modes}, their decomposition into
harmonics (photon sidebands) and their overlaps with the static eigenstates. Four
methods are available: diagonalization in the extended (Sambe) space, the one-period
propagator, a high-frequency effective Hamiltonian and a few-level rotating-wave model.
Their results agree where their regimes of validity overlap, and the manual uses these
comparisons throughout.

\section{A physics-first reading route}

Begin with the device sketches, Andreev trajectories and occupation diagrams in
\cref{ch:background}. Then use the four-state dot in \cref{ch:model} to understand
the role of charging energy. \Cref{ch:correlations} distinguishes proximity pairing
from Kondo screening, and \cref{ch:experiments} explains which transitions an
experiment can detect and how the states can be used as qubits. The first executable calculation is
in \cref{sec:first-calculation}, after the singlet and doublet states have been
defined. \Cref{ch:static} explains the approximations needed to compute them.

\section{Conventions}

% Generated by tools/gen_docs.py from docs/conventions.toml; do not edit.
\begin{table}[htb]
\centering\small
\caption{Physical and numerical conventions of \pkg{}.}\label{tab:conventions}
\begin{tabularx}{\textwidth}{>{\raggedright\arraybackslash}p{3.6cm}X}
\toprule
Quantity & Convention \\
\midrule
Units & Energies use the units of the Hamiltonian parameters, usually with $\Delta=1$; $\hbar=1$, so $\Omega$ and $1/t$ have energy units \\
Impurity & $H_d=\tfrac U2(n_d-1)^2+\epsilon(n_d-1)+B\,S^z_d$; the detuning $\epsilon=\epsilon_d+U/2$ vanishes at half filling; $B>0$ raises the spin-up level \\
Reservoirs & BCS leads with the real pairing term $-\Delta_l\sum_i(c^\dagger_{li\uparrow}c^\dagger_{li\downarrow}+\mathrm{h.c.})$, $\Delta_l>0$, and normal-state half-bandwidth $D$; normal-state density of states per spin $\rho=1/(2D)$ \\
Hybridization & $\Gamma=\sum_l\Gamma_l$ is the total hybridization, $\Gamma_l=\pi\rho\lvert t_l\rvert^2$ with $t_l$ the hopping between the dot and the contact orbital of lead $l$ \\
Phases & Factors $\mathrm{e}^{\mathrm{i}\phi_l/2}$ on the tunnelling terms $C^\dagger_{l\sigma}d_\sigma$, i.e. the physical pair potential $\Delta_l\mathrm{e}^{-\mathrm{i}\phi_l}$ and the pair amplitude $\langle c^\dagger_{li\uparrow}c^\dagger_{li\downarrow}\rangle\propto\mathrm{e}^{\mathrm{i}\phi_l}$; the reference junction uses $\phi_{L,R}=\mp\phi/2$, $\phi=\phi_R-\phi_L$ \\
Energy zero & The isolated BCS vacuum energies are subtracted: at $\Gamma=0$ the energies are the atomic energies of $H_d$, with the reservoirs in their quasiparticle vacuum \\
Current & $I=(2e/\hbar)\,\partial E/\partial\phi=(2e/\hbar)\langle\partial H/\partial\phi\rangle$, the electric current from the left to the right lead (electron charge $-e$, $e>0$); reported as $\langle\partial H/\partial\phi\rangle$ in units of $\Delta$ \\
Voltage & $V=V_L-V_R$, the electric potential of the left lead relative to the right one, so $eV=\mu_R-\mu_L$; Josephson relation $\dot\phi=2eV/\hbar$; $IV$ is the power supplied by the source \\
Sectors & Fermion parity (0 even, 1 odd), $2S_z$, optional eta parity and particle number \\
Drives & $p(t)=p_0+\sum_n A_n\cos(n\Omega t+\theta_n)$; $A$ for a monochromatic drive, $A_\phi$ (radians) for the phase and $A_\epsilon$ (energy) for the detuning \\
Harmonics & $H(t)=\sum_k H_k\mathrm{e}^{\mathrm{i}k\Omega t}$, $H_{-k}=H_k^\dagger$, time origin $t=0$ \\
Floquet states & $\psi(t)=\mathrm{e}^{-\mathrm{i}\varepsilon t}\sum_m\mathrm{e}^{\mathrm{i}m\Omega t}u_m$ with quasienergy $\varepsilon$; Sambe matrix $H_{m-m^\prime}+m\Omega\delta_{mm^\prime}$, truncated at the Floquet cutoff $\lvert m\rvert\le M$ \\
\bottomrule
\end{tabularx}
\end{table}

